% Bab III  Format Penulisan
\chapter{FORMAT PENULISAN}

Bab ini mengatur bentuk fisik dan tata tulis skripsi agar seragam. Seluruh aturan di bawah ini sudah tertanam pada template LaTeX dan Word (Proposal, Seminar Hasil, Skripsi), sehingga mahasiswa cukup mengisi teksnya.

\section{Kertas dan Margin}
Kertas yang dipergunakan adalah HVS ukuran A4 dengan berat kertas 80 g, dicetak satu muka. Naskah asli sebagai pegangan mahasiswa yang bersangkutan, sedangkan untuk Fakultas, Jurusan/Program Studi dan keperluan lainnya dapat difotokopi dan diperbanyak dengan syarat fotokopi naskah bersih dengan halaman persetujuan bertanda tangan asli.

Batas pengetikan (margin) adalah:
\begin{enumerate}
\item sisi kiri: 4 cm;
\item sisi atas: 4 cm;
\item sisi kanan: 3 cm;
\item sisi bawah: 3 cm.
\end{enumerate}
Teks, tabel, gambar, dan nomor halaman tidak boleh melewati batas margin. Nomor halaman diletakkan sesuai aturan pada subbab Penomoran Halaman.

\section{Huruf dan Spasi}
Naskah karya ilmiah diketik menggunakan komputer dengan program pengolah kata \textit{MS. Word} atau \textit{LaTeX} dengan huruf berjenis \textit{Times New Roman}. Ukuran huruf:
\begin{enumerate}
\item isi naskah, subbab, judul tabel, dan judul gambar: 12 pt;
\item judul bab (BAB I dan judulnya): 14 pt, tebal, huruf kapital;
\item catatan kecil di bawah tabel atau gambar (sumber data): 10 pt.
\end{enumerate}
Jarak antarbaris dalam teks adalah \textbf{1,5 spasi}, sedangkan abstrak, judul tabel, judul gambar, isi tabel, daftar isi, daftar pustaka, dan lampiran diketik \textbf{satu spasi}. Setelah tanda titik dan koma diberi jarak satu ketukan. Setelah tanda koma dan titik dua diberi jarak satu spasi (sebelum titik dua tidak diberi spasi), kecuali setelah tanda titik untuk kalimat baru, diberi jarak satu spasi.

\section{Alinea dan Perataan}
Pengaturan ketikan naskah dengan rata kanan-kiri (\textit{justify}). Alinea merupakan bagian dari bab atau subbab yang menceritakan suatu pernyataan tertentu dan berisi satu ide. Penulisan alinea dimulai dengan menjorok 1 cm dari batas kiri (\textit{first line indent}); jangan membuat alinea dengan menekan spasi atau \textit{tab} berulang-ulang. Antaralinea tidak diberi baris kosong tambahan.

Jika alinea yang tertulis di baris akhir di suatu halaman hanya berisi satu baris, hal itu tidak dibenarkan, oleh karena itu perlu dipindah ke halaman berikutnya. Pada akhir halaman minimal tercantum dua baris.

\section{Judul Bab dan Subbab}
Setiap bab dimulai pada \textbf{halaman baru}. Judul bab ditulis di tengah atas halaman dengan format dua baris: baris pertama “BAB” diikuti angka Romawi (BAB I), baris kedua judul bab dengan huruf kapital, seluruhnya tebal. Bagian tanpa nomor bab (KATA PENGANTAR, DAFTAR ISI, ABSTRAK, DAFTAR PUSTAKA, dan LAMPIRAN) juga ditulis tengah, kapital, tebal, dan dimulai pada halaman baru.

Subbab diketik di pinggir sisi kiri halaman dengan format berikut:
\begin{enumerate}
\item subbab tingkat 1 (misalnya 1.1): tebal, huruf kecil kecuali huruf pertama setiap kata ditulis dengan huruf kapital;
\item subbab tingkat 2 (misalnya 1.1.1): tebal, penulisan huruf sama dengan tingkat 1;
\item subbab tingkat 3 (misalnya 1.1.1.1): miring (tidak tebal), penulisan huruf sama;
\item penomoran subbab tidak melebihi tiga tingkat di bawah bab; uraian yang lebih rinci ditulis dalam bentuk butir (a, b, c).
\end{enumerate}
Jarak antara teks dengan judul subbab berikutnya adalah tiga spasi (sekitar satu baris dan setengah) dan judul subbab tidak boleh berada sendirian di bawah halaman. Satu subbab tidak boleh berdiri tunggal; bila ada 1.1 maka harus ada 1.2.

\section{Penomoran Halaman}
Susunan lembaran skripsi dimulai dari halaman sampul sampai lampiran. Penomoran halaman diatur sebagai berikut:
\begin{enumerate}
\item halaman sampul, lembar persetujuan, halaman pengesahan, dan pernyataan keaslian tidak menampilkan nomor (halaman-halaman ini juga tidak dimuat dalam daftar isi);
\item bagian awal mulai dari abstrak sampai daftar lampiran memakai angka Romawi kecil (i, ii, iii, \ldots) di tengah bawah halaman;
\item bagian isi mulai Bab I sampai lampiran memakai angka Arab (1, 2, 3, \ldots), dimulai dari angka 1 pada halaman pertama Bab I dan berlanjut sampai halaman terakhir;
\item pada halaman pertama setiap bab, nomor halaman diletakkan di \textbf{tengah bawah}; pada halaman-halaman selanjutnya, nomor halaman diletakkan di \textbf{kanan atas}.
\end{enumerate}
Nomor halaman diketik tanpa tanda baca tambahan (tanpa titik atau strip).

\section{Cetak Miring, Cetak Tebal, dan Tanda Baca}
Kata-kata bahasa di luar bahasa Indonesia baku ditulis miring, misalnya bahasa Inggris, Latin dan lain-lain (\textit{database}, \textit{server}, \textit{Dendrobium}). Istilah asing yang sudah diserap ke dalam bahasa Indonesia mengikuti ejaan baku (misalnya komputer, aplikasi, sistem). Cetak miring juga digunakan untuk judul buku, nama jurnal, dan nama prosiding dalam daftar pustaka. Cetak tebal hanya digunakan untuk judul bab dan subbab, dan untuk penekanan yang sangat perlu; hindari garis bawah.

Tanda kutip menggunakan tanda kutip ganda (“ ”). Singkatan nama dan gelar mengikuti ejaan yang berlaku. Hindari penggunaan tanda seru dan tanda tanya dalam kalimat ilmiah, kecuali pada rumusan masalah berbentuk pertanyaan.

\section{Angka, Satuan, dan Simbol}
Angka dan satuan pada awal kalimat harus ditulis dengan kata, tetapi angka yang diikuti bulan, nomor halaman, persen, waktu, dan satuan tetap ditulis dengan angka, misalnya: 1 Januari 1988; 37°C; halaman 37; 5\%; 10.30 WIB (tanda titik khusus untuk waktu) dan 8 cm. Angka bilangan bulat sampai dengan sepuluh harus ditulis dengan kata kecuali diikuti dengan satuan (kg, cm), tetapi jika lebih dari sepuluh ditulis biasa, misalnya: delapan; 2 cm; 15; 20 cm dst.

Angka yang menyatakan desimal menggunakan tanda koma (1,50 kg). Ribuan atau kelipatan ribuan ditulis dengan tanda titik setiap tiga angka, misalnya: 1.000 m; 100.000 ekor. Dalam satu tabel, jumlah digit di belakang koma harus konsisten (dengan melihat jenis variabel yang diukur) apabila dicantumkan dalam satu tabel. Satuan mengikuti Sistem Internasional (SI) dan ditulis dengan spasi setelah angka (5 V, 20 mA, 3,3 kHz). Simbol dan variabel matematika ditulis miring (\textit{x}, \textit{y}), sedangkan simbol satuan tegak.

\section{Pemisahan Kata di Akhir Baris}
Pemenggalan kata pada akhir baris harus disesuaikan dengan tata bahasa yang benar. Batas sebelah kiri dan sebelah kanan harus lurus, kata di dalam kalimat dapat direnggangkan paling banyak tiga ketukan. Kata sambung di akhir halaman diikutkan pada halaman berikutnya. Nama orang, singkatan, dan alamat URL tidak dipenggal sembarang.

\section{Tabel}
Suatu tabel harus dimuat dalam satu halaman, tidak boleh diputus atau dilanjutkan pada halaman berikutnya. Tabel yang panjang (lebih dari satu halaman) diletakkan di dalam lampiran. Ketentuan penulisan tabel:
\begin{enumerate}
\item \textbf{Nomor dan judul.} Kata “Tabel”, nomor, dan judul diletakkan di \textbf{atas} tabel, di tengah, satu spasi, dengan huruf kapital hanya pada awal kalimat dan tanpa titik di akhir judul. Kata Tabel dan nomor selalu diawali huruf kapital, dan nomor tidak boleh ditulis dengan kata.
\item \textbf{Penomoran.} Nomor tabel terdiri atas nomor bab dan nomor urut tabel di bab itu, dipisah titik. Contoh: Tabel 1.1 adalah tabel pertama pada Bab I, Tabel 3.2 adalah tabel kedua pada Bab III. Tabel pada lampiran diberi awalan nomor lampiran, misalnya Tabel L2.1.
\item \textbf{Bentuk.} Tabel dibuat \textbf{kotak penuh}: garis tunggal di atas, bawah, kiri, kanan, dan antarbaris. Judul kolom ditebalkan. Isi tabel diketik satu spasi.
\item \textbf{Sumber.} Tabel yang berasal dari data sekunder disebutkan sumbernya, ditulis satu spasi di bawah garis bawah tabel dengan huruf 10 pt.
\item \textbf{Rujukan.} Setiap tabel harus dirujuk dan dijelaskan dalam teks, misalnya “Tabel 4.1 menunjukkan …”. Tabel diletakkan sedekat mungkin setelah teks yang merujuknya.
\item Tabel merupakan hasil penelitian yang sudah diolah atau dianalisis; tabel dibuat dengan fasilitas tabel pada pengolah kata, bukan hasil tangkapan layar atau foto.
\end{enumerate}

\section{Gambar}
Gambar meliputi: grafik, diagram, monogram, foto, peta, skema, tangkapan layar, dan diagram alir. Ketentuan:
\begin{enumerate}
\item \textbf{Nomor dan judul.} Judul gambar diletakkan di \textbf{bawah} gambar, di tengah, satu spasi, ditulis “Gambar”, nomor, dan judul. Huruf selain huruf pertama dari kata pertama dalam kalimat menggunakan huruf kecil, kecuali nama dan istilah yang harus kapital.
\item \textbf{Penomoran.} Seperti tabel, nomor gambar terdiri atas nomor bab dan nomor urut gambar di bab itu. Contoh: Gambar 1.1 adalah gambar pertama pada Bab I dan Gambar 2.3 adalah gambar ketiga pada Bab II. Gambar pada lampiran diberi awalan nomor lampiran, misalnya Gambar L2.1.
\item \textbf{Kualitas.} Gambar harus jelas, resolusi memadai (minimal 300 dpi untuk cetak), dan tulisan di dalamnya terbaca pada ukuran cetak. Gambar tidak melewati margin; lebar gambar maksimal selebar teks.
\item \textbf{Warna.} Gambar boleh berwarna, tetapi harus tetap dapat dibaca bila dicetak hitam-putih; untuk grafik gunakan perbedaan pola garis atau penanda selain warna.
\item \textbf{Sumber.} Gambar yang diambil dari sumber lain wajib mencantumkan sumbernya di akhir judul gambar, misalnya “(Ghosh et al., 2010)”. Gambar karya sendiri tidak perlu sumber.
\item Setiap gambar harus dirujuk dan dijelaskan dalam teks. Gambar tidak boleh terpotong oleh batas halaman.
\end{enumerate}

\section{Persamaan}
Persamaan diketik dengan fasilitas persamaan pada pengolah kata (bukan gambar). Persamaan diletakkan di tengah baris tersendiri dan diberi nomor di sisi kanan dalam tanda kurung, mengikuti nomor bab: persamaan (2.1) adalah persamaan pertama pada Bab II. Persamaan dirujuk di dalam teks dengan menulis “Persamaan (2.1)”. Setelah persamaan, tulis “dengan nilai” atau “dengan” lalu jelaskan arti setiap simbol dan satuannya pada baris berikutnya.

\section{Kutipan dan Sitasi}
Kutipan yang ditulis pada skripsi harus sesuai dengan maksud dari tulisan pada pustaka yang diambil. Sitasi di dalam teks menggunakan gaya APA edisi ke-7 dengan format penulis--tahun dan harus sesuai dengan entri pada daftar pustaka (setiap sitasi ada di daftar pustaka, dan sebaliknya).

\subsection{Bentuk Sitasi di dalam Teks}
Penempatan nama pengarang bisa di depan, tengah atau akhir kalimat. Contoh bentuk yang lazim:
\begin{enumerate}[a.]
\item Satu penulis, di depan kalimat: Andrews (2007) menyatakan bahwa …; di akhir kalimat: … (Andrews, 2007).
\item Dua penulis: Xu dan Fischione (2012) melaporkan … ; di akhir kalimat memakai tanda \&: … (Xu \& Fischione, 2012).
\item Tiga penulis atau lebih: tulis penulis pertama diikuti \textit{et al.} untuk sumber berbahasa asing, atau \textit{dkk.} untuk sumber berbahasa Indonesia, sejak sitasi pertama: … (Capozzi et al., 2013).
\item Lembaga sebagai penulis: … (Badan Pusat Statistik, 1997).
\item Tanpa penulis: gunakan judul yang dipersingkat dicetak miring: … (\textit{Panduan penulisan karya ilmiah}, 2020).
\item Tanpa tahun: gunakan (t.t.): … (Nama Lembaga, t.t.).
\item Beberapa sumber sekaligus dipisah titik koma dan diurutkan abjad: … (AlQahtani \& Alhassany, 2013; Capozzi et al., 2013).
\item Sumber dikutip dari sumber lain (sitasi sekunder): … (Andrews, 2007, dalam Capozzi et al., 2013). Hindari bila sumber aslinya dapat diperoleh.
\end{enumerate}
Singkatan \textit{et alii} adalah \textit{et al.}, dipakai apabila pengarangnya bukan orang Indonesia. Singkatan dan kawan-kawan adalah \textit{dkk.}, dipakai apabila pengarangnya orang Indonesia. Penulisan pustaka tanpa pengarang ditulis dengan judulnya, dan apabila nama pengarang tidak ada tetapi ada lembaga yang mengeluarkan, maka ditulis nama lembaganya.

\subsection{Kutipan Langsung dan Tidak Langsung}
Kutipan tidak langsung (parafrasa) ditulis dengan kalimat sendiri tanpa tanda kutip dan wajib disertai sitasi. Kutipan langsung ditulis persis seperti aslinya dengan menyertakan nomor halaman: (Ghosh et al., 2010, hlm. 15). Kutipan langsung kurang dari 40 kata ditulis di dalam kalimat dengan tanda kutip. Kutipan langsung 40 kata atau lebih ditulis sebagai blok tersendiri, menjorok 1,27 cm dari margin kiri, satu spasi, tanpa tanda kutip, dengan sitasi di akhir. Kutipan langsung sebaiknya dibatasi karena skripsi adalah karya tulis sendiri.

\section{Singkatan, Akronim, dan Istilah Asing}
Singkatan dan akronim ditulis lengkap pada penyebutan pertama diikuti singkatannya dalam tanda kurung, misalnya \textit{Internet of Things} (IoT), selanjutnya cukup singkatannya. Singkatan yang lazim dalam bahasa Indonesia (misalnya UM Buton, S.Kom., dkk., dll.) mengikuti ejaan yang berlaku. Istilah asing yang belum ada padanan bakunya ditulis miring pada penyebutan pertama dan diberi terjemahan bila perlu.

\section{Bahasa}
Skripsi ditulis dalam bahasa Indonesia yang baik dan benar, mengikuti Ejaan Bahasa Indonesia yang berlaku, dengan ketentuan:
\begin{enumerate}
\item kalimat efektif, ringkas, dan memiliki subjek serta predikat yang jelas;
\item bentuk tidak personal, tidak menggunakan kata ganti “saya”, “kami”, atau “kita”; gunakan “penulis” atau kalimat pasif;
\item tidak menggunakan bahasa percakapan, singkatan tidak baku, atau kata-kata bermakna ganda;
\item satu alinea berisi satu ide pokok dan terdiri atas minimal tiga kalimat;
\item konsisten dalam penulisan istilah, satuan, dan nama sepanjang naskah;
\item penulisan kata “di” sebagai kata depan dipisah (di rumah), dan sebagai awalan digabung (dipakai).
\end{enumerate}

\section{Pemeriksaan Kemiripan (Plagiasi)}
Sebelum diujikan, skripsi diperiksa kemiripan tulisannya menggunakan perangkat lunak yang ditetapkan Program Studi. Batas kemiripan maksimum yang diperbolehkan ditetapkan oleh Program Studi dan hasil pemeriksaan dilampirkan pada pengajuan ujian.
